\documentclass[11pt]{article}
\usepackage[margin=1in]{geometry}
\usepackage{amsmath,amssymb,amsthm,mathtools}
\usepackage{enumitem}
\usepackage{booktabs}
\usepackage{hyperref}
\usepackage{xcolor}
\usepackage[most]{tcolorbox}

\hypersetup{colorlinks=true,linkcolor=blue!50!black,citecolor=blue!50!black,urlcolor=blue!50!black}

\newtheorem{theorem}{Theorem}[section]
\newtheorem{lemma}[theorem]{Lemma}
\newtheorem{proposition}[theorem]{Proposition}
\newtheorem{corollary}[theorem]{Corollary}
\theoremstyle{definition}
\newtheorem{definition}[theorem]{Definition}
\newtheorem{example}[theorem]{Example}
\newtheorem{remark}[theorem]{Remark}

% Status tags for reviewer checklist
\newcommand{\status}[1]{\textit{\textcolor{black!70}{[#1]}}}
\newcommand{\checkbox}{\fbox{\phantom{X}}}

\newtcolorbox{checkboxenv}[1][]{
  colback=gray!6,colframe=gray!50,boxrule=0.4pt,arc=1pt,
  left=4pt,right=4pt,top=3pt,bottom=3pt,
  title={What to check},fonttitle=\bfseries\small,#1
}

\title{\textit{p}-Adic Transfer, Signature Tails, and Homology Manifolds:\\
From a 1980 Binghamton Thesis to the 2026 Hilbert--Smith Arguments}
\author{Alan J.\ Coppola\thanks{%
Ph.D., SUNY Binghamton, 1979/80; advisor Louis~F.\ McAuley.
Email: \texttt{ajjc@optngn.com}.
Draft for private review; co-authorship invitations open.
OpenAI (2026-09-23) and A.~Dabrowski (2026-10-08) are cited as source
arguments being extended and delimited, not as co-authors of this draft.}}
\date{Draft: 18 October 2026\\[0.4em]
\small\textit{arXiv-style preprint --- not yet submitted}}

\begin{document}
\maketitle

\begin{abstract}
I revisit the Hilbert--Smith $p$-adic exclusion problem through the lens of my
1980 Binghamton dissertation on $p$-adic transfer and Smith--Yang sequences,
and through the two 2026 packaging arguments of OpenAI (Witt/signature lattice
on $S^d$) and Dabrowski (integral $L_4$ signature tails).  The classical
Yang jump $\dim_{\mathbb{Z}}(M/\mathbb{Z}_p)=n+2$ is $p$-torsion; over
$\mathbb{Q}$ the transfer is a projection and no dimension jump occurs.
I give a clean diagnosis of the break in McAuley's 2001 arXiv manuscript,
prove (relative to the 2026 manifold theorems) that connected polyhedral and
resolvable ENR homology manifolds reduce to the topological-manifold case,
and isolate the remaining crown jewel: nowhere-Euclidean Bryant--Ferry--Mio--Weinberger
ENR $\mathbb{Z}$-homology manifolds with Quinn index $I(X)\neq 1$.
There OpenAI's equal-parts lattice argument rephrases to kill every $I(X)\neq 0$,
while Dabrowski's single-$p$ cyclotomic detector fails whenever $p\mid I(X)$;
a $k$-step cyclotomic tower restores $p^k$-divisibility.
An inverse-limit trichotomy shows Raymond--Williams-style towers cannot produce
a connected locally connected ENR cohomology-manifold counterexample.
The unique remaining frontier is class HM-6
($\dim_{\mathbb{Z}}X=n<\infty$, $\dim X=\infty$), where we develop a
Borel-join Edwards--Walsh + Williams splicing program and verify its algebraic
gates computationally and in Lean~4.
Section~\ref{sec:compare} separates three layers when
$X_{\mathrm{man}}=\varnothing$: chart proofs fail, signature exclusion is
conditional on the ENR rephrase or $k$-tower, and inverse-limit
\emph{construction} is blocked---so emptiness of the manifold stratum is not
a license to build an effective $p$-adic action.  Each major statement is tagged
\textit{classical / relative / new} and carries a short reviewer checklist.
\end{abstract}

\tableofcontents
\vspace{1em}
\noindent\textit{Status of this draft.}
This is a working manuscript for reading, annotation, and possible co-authorship.
It does \textbf{not} claim an independent complete proof of the classical
Hilbert--Smith conjecture.  Theorems that depend on OpenAI or Dabrowski are
explicitly marked \status{relative}.  Lean formalization notes appear in
Appendix~\ref{app:lean}.

%======================================================================
\section{Introduction: a memoir bridge}
\label{sec:intro}

\subsection{What Hilbert--Smith asks}
A continuous effective action of a $p$-adic group $\mathbb{Z}_p$ on a connected
finite-dimensional topological manifold $M$ is conjectured to be impossible
(Hilbert's fifth problem in the compact/transformation-group form; see
\cite{newman1931,smith1941,yang1960}).  Finite cyclic groups of $p$-power order
do act freely on manifolds (lens spaces); the obstruction, if it exists, must
use the infinite descending chain of open subgroups
$p^j\mathbb{Z}_p\subset\mathbb{Z}_p$.

\subsection{What I proved in 1980}
In my dissertation~\cite{coppola1980}, advised by L.~F.\ McAuley at SUNY
Binghamton, I constructed a $p$-adic transfer homomorphism for the inverse
system of finite cyclic groups $\mathbb{Z}_{p^r}=A_p/p^r A_p$ and used it to
derive the Smith--Yang exact sequences governing $\mathbb{Z}_{p^r}$-actions on
compact metric spaces.  Two facts were already clear:
\begin{enumerate}[label=(\alph*)]
\item Over a field $L$ of characteristic $\neq p$ (in particular over
$\mathbb{Q}$), the transfer $(1/p^r)\sum_{g\in\mathbb{Z}_{p^r}}g^*$ is a
projection at every finite stage.  Passing to the limit identifies
$\check{H}^*(X/A_p;\mathbb{Q})\xrightarrow{\sim}\check{H}^*(X;\mathbb{Q})^{A_p}$.
\item Over $\mathbb{Z}$ and $\mathbb{F}_p$, the composite of transfer and
projection is multiplication by $p^r\equiv 0\pmod{p}$, producing the
Bockstein/Smith--Gysin connecting maps that raise cohomological dimension.
On a compact $n$-manifold (or suitable homology manifold) $M$, Yang's theorem
and Bredon--Raymond--Williams give
\[
\dim_{\mathbb{Z}}(M/A_p)=n+2,\qquad
H^{n+2}_c(M/A_p;\mathbb{Z})/K\cong\mathbb{Z}[1/p]/\mathbb{Z}.
\]
\end{enumerate}
Raymond and Williams~\cite{raymondwilliams1960,raymondwilliams1963} realized the
same $(n+2)$-dimensional jump on an $n$-dimensional compactum by inverse limits
of $(n+2)$-complexes with period-$p^i$ homeomorphisms retracting onto
$n$-skeleta.  Those examples are \emph{not} locally Euclidean and are not
locally connected in the cohomology-manifold sense.

\subsection{McAuley 2001 and the break}
McAuley~\cite{mcauley2001} attempted a full proof of Hilbert--Smith.
His argument uses of $M$ only local compactness, connectedness, $lc^n$, and
local homology $H_n(X,X\setminus D;\mathbb{Z}_p)\cong\mathbb{Z}_p$.
Because $\dim(M/A_p)=n+2$, saturated covers pulled back from the orbit space
have order $>n+1$.  To force unique $n$-cycles for a \v{C}ernavski\u{\i}-type
sum, he interposes a refinement $U^m$ of order $n+1$.  That cover
\textbf{cannot} be $A_p$-invariant, and the resulting chain maps fail to be
equivariant on the $n$-skeleton of the nerve.  The missing $(n+1)$-chain is
precisely where the Smith--Yang transfer cocycle lives.  I record this as
Corollary~\ref{cor:mcauley} below; it is historical diagnosis, not a
replacement proof.

\subsection{Two 2026 packagings}
In September--October 2026 two independent arguments appeared that sit in a
different family from Yang dim-raising:
\begin{itemize}
\item \textbf{OpenAI}~\cite{openai2026}: after Newman chart reduction and Haar
averaging into an odd sphere $S^d$, an integral Witt/signature lattice with
fixed denominator bound $L_d$ is divided into $p^k$ equal parts
$4/p^k$, contradicting integrality for large $k$.
\item \textbf{Dabrowski}~\cite{dabrowski2026}: an integral signature
\emph{tail} in $L_4(A_G(T))$ is $p$-divisible by Karoubi/controlled
Mayer--Vietoris, while an effective $\mathbb{Z}_p$-action realizes a class
with tail $1$ via a degree-one sphere control map and $\mathbb{CP}^2$
localization---contradicting $1\in p\mathbb{Z}$.
\end{itemize}
Both share a Euclidean shell (Newman, charts, orientation, small open
subgroups).  Their arithmetic engines differ: equal parts in a fixed lattice
versus single-$p$ cyclotomic divisibility of a tail.

\subsection{Relative-paper framing and seven contributions}
\label{sec:contributions}

\begin{checkboxenv}
\textbf{ArXiv framing (read first).}
The exclusion results Theorems~\ref{thm:E}--\ref{thm:F} are
\textbf{conditional / relative}: assuming the 2026 topological-manifold
exclusion machines of OpenAI~\cite{openai2026} or
Dabrowski~\cite{dabrowski2026}, we determine the exact boundary of $p$-adic
exclusion across the six classes HM-1--HM-6 of Definition~\ref{def:hm}.
This manuscript does \textbf{not} claim a standalone classical proof of
Hilbert--Smith on topological manifolds.  What is new is the boundary map,
the Quinn-index / $k$-tower arithmetic, the inverse-limit trichotomy, and the
HM-6 construction program.
\end{checkboxenv}

Under that framing the paper makes seven mathematical contributions:
\begin{enumerate}[label=\textbf{C\arabic*.}]
\item \textbf{Quinn-index obstruction to single-$p$ arithmetic}
(Theorem~\ref{thm:dabrowski-gap}):
on a nowhere-Euclidean BFMW ENR $\mathbb{Z}$-homology $n$-manifold
($n\ge 6$, $X_{\mathrm{man}}=\varnothing$), Dabrowski's single index-$p$ step
yields only $I(X)\in p\mathbb{Z}$, which fails to contradict whenever an odd
prime $p$ divides $I(X)$ (e.g.\ $I(X)=9$ for $p=3$).
\item \textbf{$k$-step cyclotomic tower upgrade}
(Theorem~\ref{thm:F} + Lean~4):
replacing the single index-$p$ subgroup by a $k$-step filtration forces
$p^k\mid I(X)$; choosing $p^k>|I(X)|$ kills every nonzero Quinn index.
Formalized in \texttt{CyclotomicTower.lean} / \texttt{QuinnIndex.lean}.
\item \textbf{OpenAI ENR rephrase for nowhere-Euclidean generalized manifolds}
(Theorem~\ref{thm:E}):
bypassing source-sphere compactification by keeping
$\widehat{M}=X\times S^{d-n}$ and pinching into $S^d$; equal-parts
integrality kills every $I(X)\neq 0$.
\item \textbf{Unconditional reductions for polyhedral and nonempty-stratum
homology manifolds}
(Theorems~\ref{thm:C},~\ref{thm:D}):
connected polyhedral $L$-homology manifolds have dense open
$\mathrm{Homeo}(X)$-invariant manifold strata; resolvable ENRs reduce via
$X\times\mathbb{R}^2$.
\item \textbf{One-page diagnosis of McAuley's 2001 break via the 1980 transfer}
(Corollary~\ref{cor:mcauley}):
an order-$(n+1)$ refinement cannot be $A_p$-invariant when
$\dim_{\mathbb{Z}}(M/A_p)=n+2$.
\item \textbf{Bounded-dimension inverse-limit no-go trichotomy}
(Theorem~\ref{thm:G}):
free step-kernel $C_p$ on compact bonding fibers forces Cantor fibers, an
Euler/Smith contradiction, or destruction of local $H_n$.
\item \textbf{HM-6 infinite-covering-dimension frontier}
(Theorem~\ref{thm:dranishnikov-escape},
Proposition~\ref{prop:borel-join-hm6}):
the unique class escaping OpenAI ENR rephrase, Dabrowski $k$-tower, and
bounded-dimension Euler obstruction, with an explicit Borel-join
Edwards--Walsh + Williams splicing architecture verified by a symbolic
harness (Gates G1--G4) and Lean~4 module \texttt{DranishnikovHM6.lean}.
\end{enumerate}

%======================================================================
\section{Preliminaries}
\label{sec:prelim}

\begin{definition}[Manifold stratum]
For a space $X$ and fixed $n$, write $X_{\mathrm{man}}$ for the set of points
with a neighborhood homeomorphic to $\mathbb{R}^n$.  We say $X$ is
\emph{nowhere Euclidean} if $X_{\mathrm{man}}=\varnothing$.
\end{definition}

\begin{definition}[Homology / cohomology manifolds]
An ENR $\mathbb{Z}$-homology $n$-manifold is a Euclidean neighborhood retract
$X$ with $H_*(X,X\setminus\{x\};\mathbb{Z})\cong H_*(\mathbb{R}^n,\mathbb{R}^n\setminus\{0\})$
for every $x$.  By Raymond~\cite{raymond1965}, for compact ENRs this coincides
with the ENR $\mathbb{Z}$-cohomology $n$-manifold condition.  We write
\emph{generalized $n$-manifold} in the sense of Bryant--Ferry--Mio--Weinberger
\cite{bfmw1996}.
\end{definition}

\begin{definition}[Quinn index]
For a compact connected ENR $\mathbb{Z}$-homology $n$-manifold $X$ ($n\ge 6$),
Quinn's resolution index $I(X)\in\mathbb{Z}$ satisfies
$I(X)\equiv 1\pmod{8}$.  Resolvable spaces have $I(X)=1$; BFMW construct
examples with $I(X)\neq 1$ and $X_{\mathrm{man}}=\varnothing$.
The product formula $I(X\times Y)=I(X)I(Y)$ implies that if $I(X)\neq 1$,
then $X\times Y$ is never a topological manifold for any generalized manifold
$Y$ with $I(Y)\equiv 1\pmod{8}$.
\end{definition}

\begin{definition}[Six homology-manifold classes]
\label{def:hm}
\begin{description}[style=nextline,leftmargin=1.2em]
\item[HM-1] Non-$clc$ \v{C}ech cohomology $n$-manifolds over $L$ with
$\mathrm{char}(L)\neq p$, without $\mathbb{F}_p$-regularity.
\item[HM-2] Disconnected or $0$-dimensional-junction $\mathbb{Z}$-homology
$n$-manifolds (shrinking necklace).
\item[HM-3] Connected polyhedral (simplicial/PL) $L$-homology $n$-manifolds.
\item[HM-4] Connected resolvable ENR $\mathbb{Z}$-homology $n$-manifolds
($I(X)=1$, or any connected $\mathbb{F}_p$-cm with $X_{\mathrm{man}}\neq\varnothing$).
\item[HM-5] Nowhere-Euclidean non-resolvable ENR $\mathbb{Z}$-homology
$n$-manifolds ($n\ge 6$, $I(X)\neq 1$, $\dim X=n$).
\item[HM-6] Non-ENR (Dranishnikov / Edwards--Walsh) $\mathbb{Z}$-cohomology
$n$-manifolds with \emph{finite cohomological dimension}
$\dim_{\mathbb{Z}}X=n<\infty$ and \emph{infinite topological covering dimension}
$\dim X=\infty$ ($X_{\mathrm{man}}=\varnothing$).
\end{description}
\end{definition}

%======================================================================
\section{Classical transfer and the McAuley break}
\label{sec:classical}

\begin{lemma}[Transfer over $\mathbb{Q}$]
\label{lem:A}
\status{classical --- Coppola 1980 / BRW}
Let $G$ be a compact $p$-adic group acting continuously on a compact
metrizable space $E$.  Then
\[
r^*:\check{H}^*(E/G;\mathbb{Q})\longrightarrow\check{H}^*(E;\mathbb{Q})^G
\]
is split injective.  In particular there is no rational cohomological
dimension jump.
\end{lemma}

\begin{proof}[Sketch]
At each finite stage $\mathbb{Z}_{p^r}$, the operator
$(1/p^r)\sum g^*$ is a projection onto invariants.  Inverse limits preserve
the splitting over $\mathbb{Q}$.
\end{proof}

\begin{checkboxenv}
Confirm that OpenAI's \texttt{lem:orbit-injectivity} is exactly this statement
(no covering-dimension hypothesis).  Confirm that Dabrowski never needs a
rational dimension jump.
\end{checkboxenv}

\begin{lemma}[Yang jump is $p$-torsion]
\label{lem:B}
\status{classical --- Yang / BRW / Coppola}
If $M$ is a compact connected orientable $n$-manifold (or suitable
$\mathbb{Z}$-homology $n$-manifold) with a faithful continuous
$\mathbb{Z}_p$-action, then $\dim_{\mathbb{Z}}(M/\mathbb{Z}_p)=n+2$ and the
extra cohomology is $p$-primary torsion of the form
$\mathbb{Z}[1/p]/\mathbb{Z}$ up to finite kernel.
\end{lemma}

\begin{corollary}[McAuley break]
\label{cor:mcauley}
\status{new as a one-page diagnosis; phenomenon classical}
In the setting of~\cite{mcauley2001}, any cover $U^m$ of order $n+1$ pulled
back toward $M$ cannot be $A_p$-invariant when
$\dim_{\mathbb{Z}}(M/A_p)=n+2$.  Consequently the chain-level maps used to
produce unique $n$-cycles on $N(V^m_n)$ fail to be $A_p$-equivariant; the
discrepancy is an $(n+1)$-chain carrying the Smith--Yang transfer cocycle.
\end{corollary}

\begin{checkboxenv}
Read McAuley Lemmas 6--8 (\texttt{arXiv:math/0103215v8}, pp.~near 33 and
1460--1488 in the PDF text layer).  Check that order$(U^m)=n+1$ is forced by
wanting unique $n$-cycles, and that this contradicts $A_p$-invariance when the
orbit space has dimension $n+2$.
\end{checkboxenv}

%======================================================================
\section{Counterexamples on weakened hypotheses}
\label{sec:counterex}

\begin{proposition}[Necklace counterexample]
\label{prop:necklace}
\status{new as an explicit HM-2 class; ingredients classical}
Fix $m\ge 1$.  On the $i$-th copy $M_i=S^{2m-1}$, let $\mathbb{Z}_p$ act via
$\mathbb{Z}_p\twoheadrightarrow C_{p^i}\subset S^1$.  Embed the spheres as
radius $2^{-i}$ and shrink to a junction point $x_\infty$.  The resulting
compact connected metrizable space $X$ carries a faithful continuous
$\mathbb{Z}_p$-action, but local homology at $x_\infty$ has infinite rank, so
$X$ is not a $\mathbb{Z}$-homology manifold at the junction.
\end{proposition}

\begin{corollary}[Coefficient mismatch]
\label{cor:hm1}
\status{classical mechanism --- Raymond 1961/1965}
If one drops $lc^n$ over $\mathbb{Z}$ and $\mathbb{F}_p$-cohomology-manifold
regularity, Vietoris maps over $\mathbb{Q}$ may have Cantor $\mathbb{Z}_p$-fibers
because $\check{H}^*(\mathbb{Z}_p;\mathbb{Q})$ is trivial in positive degrees.
Neither 2026 proof leaks onto this class: both require local connectedness and
Newman-type rigidity over $\mathbb{F}_p$.
\end{corollary}

%======================================================================
\section{Polyhedral and resolvable reductions}
\label{sec:reductions}

\begin{theorem}[Polyhedral reduction]
\label{thm:C}
\status{new packaging; classical PL local duality}
Let $L\neq 0$ be a commutative ring with unit and let $X=|K|$ be a connected
second-countable polyhedral $L$-homology $n$-manifold ($n\ge 1$).  Then:
\begin{enumerate}[label=(\roman*)]
\item $X_{\mathrm{man}}$ contains $X\setminus|K^{(n-2)}|$ and is a connected,
dense, open, $\mathrm{Homeo}(X)$-invariant topological $n$-manifold without
boundary;
\item every faithful continuous $\mathbb{Z}_p$-action on $X$ restricts to a
faithful continuous action on $X_{\mathrm{man}}$;
\item consequently both OpenAI~\cite{openai2026} and Dabrowski~\cite{dabrowski2026},
if accepted for topological manifolds, exclude faithful continuous
$\mathbb{Z}_p$-actions on $X$.
\end{enumerate}
\end{theorem}

\begin{proof}
In a simplicial $L$-homology $n$-manifold the link of every $(n-1)$-simplex
has $\tilde{H}_0\cong L$, hence exactly two vertices.  Open stars off the
$(n-2)$-skeleton are Euclidean.  Connectivity of the dual
$(n,n-1)$-gallery graph prevents $|K^{(n-2)}|$ from separating $X$.
Homeomorphisms preserve $X_{\mathrm{man}}$; density plus Hausdorffness makes
restriction $\mathrm{Homeo}(X)\to\mathrm{Homeo}(X_{\mathrm{man}})$ injective.
Apply the topological-manifold theorem on $X_{\mathrm{man}}$.
\end{proof}

\begin{remark}[Direct Dabrowski route on polyhedral $\mathbb{Q}$-hm]
Even without deleting $|K^{(n-2)}|$, Dabrowski's dual-block carrier acyclicity
(his Lemma~7.4) holds verbatim on any polyhedral $\mathbb{Q}$-homology
manifold: only rational link acyclicity is used.
\end{remark}

\begin{checkboxenv}
Verify gallery connectivity for connected polyhedral homology manifolds over
an arbitrary coefficient ring $L\neq 0$.  Confirm that Dabrowski Lemma~7.4
really uses only $\mathbb{Q}$-link acyclicity.
\end{checkboxenv}

\begin{theorem}[Resolvable ENR / nonempty manifold stratum]
\label{thm:D}
\status{relative to Newman on $\mathbb{F}_p$-cm + Edwards--Quinn}
Let $X$ be a connected Hausdorff second-countable $\mathbb{F}_p$-cohomology
$n$-manifold of finite covering dimension.
\begin{enumerate}[label=(\roman*)]
\item If $X_{\mathrm{man}}\neq\varnothing$, then any faithful continuous
$\mathbb{Z}_p$-action induces a faithful action of some open subgroup
$p^j\mathbb{Z}_p\cong\mathbb{Z}_p$ on a Euclidean open set
$O\subset X_{\mathrm{man}}$, so both 2026 arguments apply as written.
\item If $I(X)=1$ (resolvable case), then $X\times\mathbb{R}^2$ is a
topological $(n+2)$-manifold by Daverman and Edwards--Quinn, and a faithful
action on $X$ extends to $X\times\mathbb{R}^2$.
\end{enumerate}
\end{theorem}

\begin{proof}[Sketch of (i)]
Newman's theorem for connected $\mathbb{F}_p$-cohomology manifolds of finite
covering dimension~\cite{smith1941,connerfloyd1959} implies that no nontrivial
$p$-power periodic homeomorphism fixes a nonempty open set.  The OpenAI
chart-reduction argument then produces a faithful small-subgroup action on a
coordinate ball.
\end{proof}

\begin{checkboxenv}
Confirm Newman for connected $\mathbb{F}_p$-cm (Conner--Floyd Cor.~6.2 / Bredon).
Confirm Daverman: $I(X)=1\Rightarrow X\times\mathbb{R}^2$ has DDP and is a
manifold by Edwards--Quinn.
\end{checkboxenv}

%======================================================================
\section{The crown jewel: nowhere-Euclidean ENR homology manifolds}
\label{sec:hm5}

By Theorem~\ref{thm:D}, the only connected $\mathbb{Z}$-homology $n$-manifolds
not immediately reduced to a Euclidean chart are those with
$X_{\mathrm{man}}=\varnothing$.  Such spaces exist for every $n\ge 6$
\cite{bfmw1996}, with arbitrary Quinn index $I(X)\equiv 1\pmod{8}$,
$I(X)\neq 1$.

\begin{theorem}[Dabrowski arithmetic fails when $p\mid I(X)$]
\label{thm:dabrowski-gap}
\status{new as a delimitation of Dabrowski}
Suppose one replaces the compactified coordinate sphere in Dabrowski
\S\S8--13 by a compact connected orientable non-resolvable ENR
$\mathbb{Z}$-homology $n$-manifold $X$ with Quinn index $I(X)$.  After
Karoubi localization cutting $X\times\mathbb{CP}^2$ down to a scalar germ,
the signature tail is
\[
\sigma_{\mathrm{tail}}(\alpha)=\bigl[(I(X),I(X),I(X),\ldots)\bigr].
\]
Dabrowski's divisibility uses only $\Phi_{p^{a_i}}(1)=p$, hence only
single-$p$ divisibility $\sigma_{\mathrm{tail}}\in p\mathbb{Z}^{\mathbb{N}}/\mathbb{Z}^{(\mathbb{N})}$.
For every odd prime $p$ dividing $I(X)$ (e.g.\ $I(X)=9$, $p=3$), the
relation $I(X)\in p\mathbb{Z}$ is true and yields \textbf{no contradiction}.
\end{theorem}

\begin{theorem}[OpenAI-style ENR rephrase]
\label{thm:E}
\status{relative to OpenAI; new as an ENR packaging}
Let $X$ be a compact connected orientable ENR $\mathbb{Z}$-homology
$n$-manifold with a continuous $\mathbb{Z}_p$-action.  Rephrase
OpenAI~\cite{openai2026} by setting $\widehat{M}=X\times S^{d-n}$
($d>n$ odd, $d\ge 3$) as \emph{source}, keeping the \emph{target} test space
the standard sphere $S^d$.  Then:
\begin{enumerate}[label=(\roman*)]
\item Newman holds on $X$ as a connected $\mathbb{F}_p$-cm;
\item $\widehat{M}$ is a retract of a compact finite polyhedron, so derived
pushforwards land in the idempotent-complete test category
$\mathcal{T}(S^d)$ without triangulating $\widehat{M}$;
\item moving the test uses Poincar\'e--Alexander duality on the homology
manifold $W\subset\widehat{M}$ rather than sliding Euclidean balls;
\item the denominator lattice $L_d^{-1}\mathbb{Z}\overline{u}_d$ is a
property of the \emph{target} $S^d$ alone, so the equal-parts identity
\[
\overline{z}_i(a)=\frac{4\,c_X}{p^k}\,\overline{u}_d
\]
contradicts integrality once $p^k>4|c_X|L_d$, for \textbf{every} prime $p$
and every $c_X=I(X)\neq 0$.
\end{enumerate}
\end{theorem}

\begin{checkboxenv}
\textbf{Critical relative steps.}
(1)~Idempotent-completeness of $\mathcal{T}(S^d)$ absorbs ENR retracts.
(2)~Alexander duality $\check{H}^d(W^+\setminus q^{-1}(V);\mathbb{Q})=0$
replaces the sliding-ball homotopy.
(3)~The lattice bound $L_d$ depends only on the target sphere.
Mark WEAK if any of these is not accepted as written in OpenAI.
\end{checkboxenv}

\begin{theorem}[Dabrowski $k$-tower upgrade]
\label{thm:F}
\status{new arithmetic proposal; relative to Dabrowski geometric engine}
Replace the single index-$p$ subgroup detector in Dabrowski~\S8 by a
$k$-step filtration
\[
H_j\supset H_{j+1}\supset\cdots\supset H_{j+k}
\]
with $k$ equivariant sphere detectors for the free $C_p$-quotients
$H_{j+r-1}/H_{j+r}$.  The $k$ commuting nilpotent operators
$\nu_r=p\,\mathrm{id}-\tau\otimes(-)$ force the integer multiplicity
polynomial $F(t)\in\mathbb{Z}[t]$ to be divisible by
$\prod_{r=0}^{k-1}\Phi_{p^{a_i-r}}(t)$, hence
\[
p^k=\prod_{r=0}^{k-1}\Phi_{p^{a_i-r}}(1)\ \Big|\ F(1)=\sigma(V_i,b_i).
\]
Choosing $k$ with $p^k>|I(X)|$ upgrades the contradiction from
$I(X)\in p\mathbb{Z}$ to $I(X)\in p^k\mathbb{Z}$, killing every Quinn index
compatible with BFMW (combined with $\varepsilon$-Poincar\'e polyhedral
approximations in place of a fixed triangulation).
\end{theorem}

\begin{checkboxenv}
Check that $k$ commuting nilpotents really multiply cyclotomic factors in
$\mathbb{Z}[t]$ without introducing denominators (the polynomial divisibility
step is formalized and compiled in \texttt{CoppolaHS/CyclotomicTower.lean}).
Check that BFMW $\varepsilon$-approximations supply the controlled input
Dabrowski needs in place of a single triangulation.  This theorem is the
natural invitation point for Dabrowski as co-author or correspondent.
\end{checkboxenv}

\subsection{Plain-English guide to Theorems~\ref{thm:dabrowski-gap}--\ref{thm:F}}
\label{sec:plain-english}

Because the notation above mixes controlled $L$-theory, derived sheaf
categories, and Quinn's surgery index, it is worth stating in plain English
what Theorems~\ref{thm:dabrowski-gap},~\ref{thm:E}, and~\ref{thm:F} actually
say and why Theorems~\ref{thm:E} and~\ref{thm:F} are marked \status{relative}.

\begin{enumerate}[label=\textbf{\arabic*.}]
\item \textbf{Are Theorems~\ref{thm:E} and~\ref{thm:F} genuine theorems?}
Yes---they are \emph{conditional (relative) theorems}.  The tag
\status{relative} means that this short paper does not re-prove the full
2026 OpenAI or Dabrowski manifold machines from scratch; instead it proves that
\emph{if} the OpenAI sheaf-Witt machine on $S^d$ holds (or \emph{if}
Dabrowski's controlled $L$-theory machine holds on $\varepsilon$-Poincar\'e
spaces), then the exclusion of effective $\mathbb{Z}_p$-actions extends from
topological manifolds to every compact connected orientable ENR
$\mathbb{Z}$-homology manifold, even when $X_{\mathrm{man}}=\varnothing$.
\item \textbf{What is the ``single-$p$ gap'' in
Theorem~\ref{thm:dabrowski-gap}?}
Dabrowski's proof on a manifold extracts an integer signature $\sigma$ and
proves two things:
\begin{enumerate}[label=(\alph*)]
\item algebraically, using \emph{one} step $H_j/H_{j+1}\cong C_p$ of the
$p$-adic group and the cyclotomic identity $\Phi_{p^a}(1)=p$, the integer
$\sigma$ must be divisible by $p$ ($\sigma\in p\mathbb{Z}$);
\item geometrically, on a Euclidean manifold chart, $\sigma=1$, contradicting
$1\in p\mathbb{Z}$.
\end{enumerate}
On a nowhere-Euclidean BFMW homology manifold $X$ ($X_{\mathrm{man}}=\varnothing$),
the local signature is no longer $1$: it is \emph{Quinn's resolution index}
$I(X)\equiv 1\pmod{8}$ (for instance $I(X)=9$, $25$, or $-7$).  If $I(X)=9$
and $p=3$, then Dabrowski's algebraic step says ``$9$ is divisible by $3$,''
which is \emph{true} ($9=3\times 3$) and gives \textbf{no contradiction}.
\item \textbf{What does ``OpenAI ENR rephrase'' (Theorem~\ref{thm:E}) mean?}
In OpenAI's manifold proof, one picks a Euclidean coordinate chart
$\mathbb{R}^n\subset M$, compactifies it into a \emph{source} sphere $S^d$,
and maps that source sphere to a \emph{target} test sphere $S^d$ through the
orbit space.  When $X_{\mathrm{man}}=\varnothing$, $X$ has no Euclidean chart,
so there is no source sphere.  The \emph{ENR rephrase} is a simple geometric
bypass: do \emph{not} replace the source by a sphere.  Keep
$\widehat{M}=X\times S^{d-n}$ itself as the source space, and pinch a small
neighborhood into the standard target sphere $S^d$.  Because $X$ is a compact
\emph{ENR} (a retract of a finite polyhedron $B$), the pushed-forward sheaf
complex from $\widehat{M}$ is a direct summand of the pushforward from $B$, so
it automatically lies in OpenAI's idempotent-complete test category
$\mathcal{T}(S^d)$ on the target sphere without needing a chart or a
triangulation.  Instead of dividing $4$ into $p^k$ equal parts, one divides
$4\,I(X)$ into $p^k$ equal parts in a fixed denominator lattice
$L_d^{-1}\mathbb{Z}$.  Because $k$ can be chosen arbitrarily large
($p^k>4|I(X)|L_d$), this kills \emph{every} $I(X)\neq 0$, even when
$p\mid I(X)$.
\item \textbf{What does ``Dabrowski $k$-tower upgrade'' (Theorem~\ref{thm:F})
mean?}
Instead of using a single subgroup step $H_j\supset H_{j+1}$ (which only divides
by a single factor of $p$), use a \emph{tower of $k$ nested steps}
\[
H_j\supset H_{j+1}\supset H_{j+2}\supset\cdots\supset H_{j+k}
\]
inside $\mathbb{Z}_p$.  Each of the $k$ steps contributes a distinct
cyclotomic polynomial $\Phi_{p^{a-r}}(t)$ dividing the character polynomial
$F(t)\in\mathbb{Z}[t]$.  Evaluating at $t=1$ multiplies the $k$ factors
$\Phi_{p^{a-r}}(1)=p$, forcing $\sigma=I(X)$ to be divisible by $p^k$
(compiled in Lean~4 in \texttt{CoppolaHS/CyclotomicTower.lean}).  When
$I(X)=9$ and $p=3$, taking $k=3$ forces $27\mid 9$, restoring the
contradiction.
\end{enumerate}

%======================================================================
\section{Settled exclusion vs open construction when $X_{\mathrm{man}}=\varnothing$}
\label{sec:compare}

The crown-jewel Theorems~\ref{thm:dabrowski-gap}--\ref{thm:F} and the
inverse-limit Theorem~\ref{thm:G} below are easy to misread as saying
``every homology manifold forbids faithful $\mathbb{Z}_p$-actions'' or,
conversely, ``when $X_{\mathrm{man}}=\varnothing$ a counterexample is still
available.''  Neither reading is correct.  This section records the logical
status explicitly.

\subsection{The question}
\label{sec:settled-question}

\begin{enumerate}[label=\textbf{Q\arabic*.}]
\item Do \emph{all} homology / cohomology manifolds exclude effective continuous
$\mathbb{Z}_p$-actions?
\item If not, can one still \emph{construct} a faithful continuous
$\mathbb{Z}_p$-action on a connected ENR $\mathbb{Z}$-homology manifold with
$X_{\mathrm{man}}=\varnothing$?
\end{enumerate}

\begin{proposition}[Short answers]
\label{prop:short-answers}
\status{new packaging of prior theorems}
\begin{enumerate}[label=(\roman*)]
\item \textbf{No to Q1.}  Weakened classes admit actions: HM-1 (coefficient /
$lc^n$ mismatch) and HM-2 (shrinking necklace) are counterexample classes
(Section~\ref{sec:counterex}).  The 2026 arguments never claimed those scopes.
\item \textbf{For connected ENR $\mathbb{Z}$-homology manifolds the answer
splits on $X_{\mathrm{man}}$ and $I(X)$:}
\begin{itemize}
\item if $X_{\mathrm{man}}\neq\varnothing$, or $X$ is polyhedral, or
$I(X)=1$, then faithful continuous $\mathbb{Z}_p$-actions are excluded
\emph{relative to} the topological-manifold Hilbert--Smith theorems
(Theorems~\ref{thm:C},~\ref{thm:D});
\item if $X_{\mathrm{man}}=\varnothing$ and $I(X)\neq 1$ (BFMW), then
\emph{as-written} OpenAI/Dabrowski chart arguments do not apply, but
\begin{itemize}
\item accepting the OpenAI ENR rephrase (Theorem~\ref{thm:E}), or the
Dabrowski $k$-tower (Theorem~\ref{thm:F}), again excludes every faithful
continuous action;
\item independently, Theorem~\ref{thm:G} blocks the standard
\emph{constructive} route (inverse limits of simplicial $C_{p^i}$-actions
with free step kernels).
\end{itemize}
\end{itemize}
\item \textbf{What remains open} is only a narrow logical gap: a faithful
continuous $\mathbb{Z}_p$-action on some BFMW space with
$X_{\mathrm{man}}=\varnothing$ and $I(X)\neq 1$, built by a method that is
\emph{neither} an inverse-limit tower of the Raymond--Williams type
\emph{nor} subject to the geometric inputs of Theorems~\ref{thm:E}
or~\ref{thm:F}.  No such construction is known, and Theorem~\ref{thm:G}
explains why the usual attempt fails.
\end{enumerate}
\end{proposition}

\subsection{Case $X_{\mathrm{man}}\neq\varnothing$: exclusion is immediate}
\label{sec:settled-nonempty}

\begin{theorem}[Chart reduction when a manifold point exists]
\label{thm:settled-chart}
\status{restatement of Theorem~\ref{thm:D}(i); classical Newman}
Let $X$ be a connected Hausdorff second-countable $\mathbb{F}_p$-cohomology
$n$-manifold of finite covering dimension.  If
$X_{\mathrm{man}}\neq\varnothing$, then a faithful continuous
$\mathbb{Z}_p$-action on $X$ induces a faithful continuous action of an open
subgroup $p^j\mathbb{Z}_p\cong\mathbb{Z}_p$ on a Euclidean open set
$O\subset X_{\mathrm{man}}$.  In particular, there is \textbf{no} remaining
construction problem in this case: any purported action is already an action
on a topological manifold chart, excluded by the 2026 manifold theorems (or by
classical Hilbert--Smith if those are granted).
\end{theorem}

Polyhedral spaces (Theorem~\ref{thm:C}) always have dense
$X_{\mathrm{man}}$.  Resolvable ENRs with $I(X)=1$ reduce by
$X\times\mathbb{R}^2$ (Theorem~\ref{thm:D}(ii)) even if the Disjoint Disks
Property fails on $X$ itself.

\subsection{Case $X_{\mathrm{man}}=\varnothing$: three layers, not one}
\label{sec:settled-empty}

Nowhere-Euclidean BFMW spaces exist for $n\ge 6$ with $I(X)\equiv 1\pmod{8}$
and $I(X)\neq 1$ \cite{bfmw1996}.  For those spaces the situation has
\emph{three} distinct layers.

\paragraph{Layer 1 --- As-written 2026 proofs do not fire.}
Both OpenAI and Dabrowski begin by choosing a Euclidean chart (or a free point
with a contractible Euclidean neighborhood).  If
$X_{\mathrm{man}}=\varnothing$, that first step is unavailable.
Theorem~\ref{thm:dabrowski-gap} records a further arithmetic failure mode for
\emph{naive} Dabrowski transplantation: single-$p$ cyclotomic divisibility does
not contradict $I(X)$ when $p\mid I(X)$.  This is a delimitation, not a
construction.

\paragraph{Layer 2 --- Conditional exclusion (Theorems~\ref{thm:E},~\ref{thm:F}).}
\begin{itemize}
\item \textbf{OpenAI route.}  Keep the \emph{target} sphere $S^d$ and replace
the source by $\widehat{M}=X\times S^{d-n}$.  Newman still holds on the
connected $\mathbb{F}_p$-cm $X$; ENR retracts land in the idempotent-complete
test category; the equal-parts lattice contradiction kills every
$I(X)\neq 0$, including those divisible by $p$.
\item \textbf{Dabrowski route.}  Upgrade the single index-$p$ detector to a
$k$-step cyclotomic tower so that $p^k\mid\sigma$ with $p^k>|I(X)|$.
\end{itemize}
If either geometric package is accepted, then \textbf{no} faithful continuous
$\mathbb{Z}_p$-action exists on any compact connected orientable ENR
$\mathbb{Z}$-homology manifold---including those with
$X_{\mathrm{man}}=\varnothing$.

\paragraph{Layer 3 --- Construction is blocked for inverse limits
(Theorem~\ref{thm:G}).}
Even if one temporarily suspends Theorems~\ref{thm:E} and~\ref{thm:F}, the
standard way one might try to \emph{build} an action---compatible simplicial
$C_{p^i}$-actions on an inverse system $E_i$, as in Raymond--Williams---cannot
produce a connected locally connected ENR cohomology manifold.  Free step
kernels on bonding fibers force Cantor failure of $lc^0$, Euler/Smith
obstructions to acyclicity, or non-acyclic splicing that destroys local
$\mathbb{Z}$-homology.  The only inverse-limit escape is the HM-2 necklace,
which is not a homology manifold at the junction.

\begin{theorem}[No inverse-limit ENR counterexample when $X_{\mathrm{man}}=\varnothing$]
\label{thm:no-ilim-enr}
\status{corollary of Theorem~\ref{thm:G}; clarified for constructors}
There is no inverse system of finite $n$-dimensional simplicial complexes with
compatible $C_{p^i}$-actions and free step-kernel action on bonding fibers
whose inverse limit is a connected locally connected ENR
$\mathbb{Z}$-cohomology $n$-manifold (resolvable or not).  In particular, one
cannot construct a BFMW counterexample to $p$-adic exclusion by a
Raymond--Williams-style tower.
\end{theorem}

\subsection{What a would-be constructor must achieve}
\label{sec:constructor-burden}

To exhibit a faithful continuous $\mathbb{Z}_p$-action on a space $X$ with
$X_{\mathrm{man}}=\varnothing$ that escapes this note, one would need all of
the following simultaneously:
\begin{enumerate}[label=(\alph*)]
\item $X$ is a connected ENR $\mathbb{Z}$-homology $n$-manifold with
$I(X)\neq 1$ (so $X\times Y$ never becomes a manifold);
\item the action is jointly continuous and faithful;
\item the geometric inputs of Theorem~\ref{thm:E} fail (e.g.\ the pushed class
is not equal-parts-integrable in $\mathcal{T}(S^d)$, or Newman/averaging
cannot be run on $X$), \emph{and} the $k$-tower inputs of
Theorem~\ref{thm:F} fail;
\item the action is not realized by an inverse limit of the type forbidden by
Theorem~\ref{thm:G}.
\end{enumerate}
Items (a)--(b) are the BFMW wish-list; (c)--(d) are the actual burden.
Until such an example appears, the correct summary is:

\begin{quote}
\itshape
Within the finite-dimensional (ENR) class, homology manifolds with
$X_{\mathrm{man}}\neq\varnothing$, polyhedral triangulations, or $I(X)=1$ are
excluded via chart reduction, and nowhere-Euclidean BFMW spaces ($I(X)\neq 1$)
are excluded relative to Theorem~\ref{thm:E} (OpenAI ENR rephrase) or
Theorem~\ref{thm:F} (Dabrowski $k$-tower), while fixed-dimension simplicial
inverse limits are blocked by Theorem~\ref{thm:G}.
To search for a \textbf{non-necklace inverse-limit example} of a
$\mathbb{Z}_p$-action on a homology manifold, one must drop finite
\textbf{topological} dimension ($\dim X<\infty$, hence ENR) while keeping finite
\textbf{cohomological} dimension ($\dim_{\mathbb{Z}}X=n<\infty$), as analyzed in
Section~\ref{sec:nogo} below.
\end{quote}

\begin{checkboxenv}
Verify that all three exclusions for $X_{\mathrm{man}}=\varnothing$
(Theorem~\ref{thm:E}, Theorem~\ref{thm:F}, and Theorem~\ref{thm:G}) depend
crucially on \emph{finite topological covering dimension} ($\dim X<\infty$ /
ENR / bounded-dimension complexes $E_i$), and that dropping $\dim X<\infty$ while
preserving $\dim_{\mathbb{Z}}X=n$ is the unique way to reopen inverse-limit
constructions without necklace pinch points.
\end{checkboxenv}

\subsection{OpenAI vs Dabrowski packaging (reminder)}
\label{sec:packaging-table}

\begin{center}
\begin{tabular}{@{}llp{5.0cm}@{}}
\toprule
Ingredient & OpenAI & Dabrowski \\
\midrule
Euclidean shell & Newman + chart + averaging & Newman + free point + sphere control \\
Arithmetic object & Witt/signature lattice on $S^d$ & $L_4$ integer signature tail \\
Divisibility & equal parts $4/p^k$, arbitrary $k$ & single $p$ via $\Phi_{p^a}(1)=p$ \\
Realization & degree-1 pinch through orbit space & $\mathbb{CP}^2$ germ after localization \\
ENR $I(X)\neq 1$ & rephrases (Thm~\ref{thm:E}) & needs $k$-tower (Thm~\ref{thm:F}) \\
Lean & none for this proof & large \texttt{HSFormal} tree \\
\bottomrule
\end{tabular}
\end{center}

Shared Euclidean ingredients: Newman rigidity, local charts/orientation,
arbitrarily small open subgroups of $\mathrm{Homeo}$, and an integral
signature lattice (Witt or $L$-theory).  Sponges/solenoids/Menger fail the
Euclidean shell, not the arithmetic---so the obstruction, if correct, is
Euclidean rather than purely dynamical.  When $X_{\mathrm{man}}=\varnothing$
the shell must be rephrased (Theorem~\ref{thm:E}) or arithmetically upgraded
(Theorem~\ref{thm:F}); it is not simply abandoned.

%======================================================================
\section{Inverse-limit no-go in finite dimension and the $\dim_{\mathbb{Z}}<\infty$, $\dim=\infty$ frontier}
\label{sec:nogo}

\subsection{The bounded-dimension trichotomy}

\begin{theorem}[Inverse-limit trichotomy in bounded topological dimension]
\label{thm:G}
\status{new}
Suppose $X=\varprojlim(E_i,\pi_i)$ is an inverse limit of finite
$n$-dimensional simplicial complexes with compatible simplicial actions of
$C_{p^i}=\mathbb{Z}_p/p^i\mathbb{Z}_p$, as in Raymond--Williams.  Let
$K_i=\ker(C_{p^i}\to C_{p^{i-1}})\cong C_p$ act on a bonding fiber
$F_x=\pi_i^{-1}(x)$.  At a regular free orbit of $C_{p^i}$ on $E_i$, the
step kernel $K_i$ acts freely on the compact fiber $F_x$, forcing one of:
\begin{enumerate}[label=\textbf{Case \arabic*:}]
\item \textbf{Zero-dimensional fibers} ($|F_x|=p$).
Then $X$ accumulates Cantor $\mathbb{Z}_p$-fibers, fails $lc^0$, and is never
a cohomology manifold over any ring.
\item \textbf{Cell-like / $\mathbb{Q}$-acyclic connected fibers at stage $i$.}
$\mathbb{F}_p$-acyclic free $C_p$-actions on finite-dimensional compacta are
impossible by Smith.  $\mathbb{Q}$-acyclic finite free $C_p$-complexes are
impossible by Euler characteristic:
$\chi(F_x)=p\cdot\chi(F_x/C_p)\equiv 0\pmod{p}$, while acyclicity forces
$\chi=1$.
\item \textbf{Connected fibers with $\chi\equiv 0\pmod{p}$ and nontrivial
homology in degrees $\le n$.}
Splicing non-acyclic fibers of dimension $\le n$ (Raymond--Williams blocks,
Menger $\mu^n$) destroys local homology
$H_n(X,X\setminus\{x\})\cong\mathbb{Z}$ along the splicing locus.
\end{enumerate}
Thus no bounded-dimension tower yields a connected locally connected ENR
cohomology manifold with a free step-kernel action on bonding fibers.
\end{theorem}

\begin{corollary}[Only necklace escape in bounded topological dimension]
\label{cor:H}
\status{new}
Within bounded-dimension simplicial towers, the only escape that preserves open
manifold strata is to allow fixed points of $K_i$ on bonding necks---i.e., a
finite necklace/wedge of manifolds or lens spaces joining along singular
points, converging to the HM-2 shrinking-necklace class of
Proposition~\ref{prop:necklace}, where local homology fails at the accumulation
point.
\end{corollary}

\begin{checkboxenv}
Case~2 Euler obstruction: confirm $\chi=1$ for $\mathbb{Q}$-acyclic finite
complexes and $\chi\equiv 0\pmod{p}$ under free $C_p$.  Case~3: confirm that
bounded-dimension RW/Menger splicing kills local $\mathbb{Z}$-homology.
\end{checkboxenv}

\subsection{How to escape the necklace trap: $\dim_{\mathbb{Z}}X=n<\infty$ with $\dim X=\infty$ (HM-6)}
\label{sec:dranishnikov}

If one wants to construct a \textbf{non-necklace} inverse-limit example of a
faithful $\mathbb{Z}_p$-action on a connected $\mathbb{Z}$-homology (or
\v{C}ech $\mathbb{Z}$-cohomology) $n$-manifold $X=\varprojlim(E_i,\pi_i)$, the
analysis above pinpoints the exact hypothesis that must be relaxed:
\emph{finite topological covering dimension} ($\dim X<\infty$).

Recall the fundamental dichotomy of cohomological dimension theory due to
Edwards--Walsh~\cite{walsh1981}, Dranishnikov~\cite{dranishnikov1988},
Mitchell~\cite{mitchell1990}, and
Dranishnikov--Ferry--Weinberger~\cite{dfw2006} and
Dydak--Walsh~\cite{dydakwalsh1993}:
\begin{itemize}
\item A compact metric $\mathbb{Z}$-cohomology $n$-manifold ($clc_{\mathbb{Z}}$,
with $H^*(X,X\setminus\{x\};\mathbb{Z})\cong H^*(\mathbb{R}^n,\mathbb{R}^n\setminus\{0\};\mathbb{Z})$
at every point) is an \textbf{ENR} (equivalently, an \textbf{ANR})
\emph{if and only if} its topological covering dimension is finite
($\dim X=n$).
\item Using Dranishnikov's solution to Alexandroff's problem ($\dim_{\mathbb{Z}}X=3$,
$\dim X=\infty$)~\cite{dranishnikov1988} and Edwards--Walsh cell-like
resolutions, there exist compact connected $\mathbb{Z}$-cohomology
$n$-manifolds (\emph{class HM-6}) that are cell-like images of $n$-manifolds or
generalized $n$-manifolds and satisfy
\[
\dim_{\mathbb{Z}}X = n < \infty, \qquad \dim X = \infty, \qquad
X_{\mathrm{man}}=\varnothing.
\]
Such spaces are inverse limits of finite polyhedra $E_i$ whose topological
dimensions \emph{grow without bound} ($\dim E_i = d_i \to \infty$), while the
bonding maps $\pi_{i+1}:E_{i+1}\to E_i$ kill all cohomology in degrees $>n$
via Eilenberg--MacLane $K(\mathbb{Z},n)$ / Postnikov resolution blocks.
\end{itemize}

\begin{theorem}[Why $\dim_{\mathbb{Z}}X=n<\infty$, $\dim X=\infty$ reopens all three doors]
\label{thm:dranishnikov-escape}
\status{new}
Let $X=\varprojlim(E_i,\pi_i)$ be a compact connected $\mathbb{Z}$-cohomology
$n$-manifold with $\dim_{\mathbb{Z}}X=n<\infty$ and $\dim X=\infty$ (so $X$ is
not an ENR and $\dim E_i = d_i \to \infty$).  Then:
\begin{enumerate}[label=(\roman*)]
\item \textbf{Theorem~\ref{thm:E} (OpenAI ENR rephrase) does not apply:}
because $\dim X=\infty$, $X$ cannot be embedded in any Euclidean space
$\mathbb{R}^N$ and is not a retract of any finite polyhedron $B$.  Hence
$Rf_*\underline{\mathbb{C}}_{\widehat{M}}$ is not a direct summand of a
polyhedral generator in $\mathcal{T}_0(S^d)$, and Verdier biduality
$a_X^!\mathbb{R}\simeq\underline{\mathbb{R}}_X[n]$ loses its finite $c$-soft
dimension hypothesis.
\item \textbf{Theorem~\ref{thm:F} (Dabrowski $k$-tower) does not apply:}
because $\dim X=\infty$, $X$ is not an ANR/ENR and does not admit
bounded-dimension controlled polyhedral Poincar\'e approximations.
\item \textbf{Case~2 of Theorem~\ref{thm:G} (single-stage Euler obstruction) is
bypassed:} when the fiber dimension $\dim F_x = d_i > n$ is allowed to grow
with $i$, a finite complex $F_x$ with a free $C_p$-action can be
\textbf{$\mathbb{Z}$-acyclic through degree $n+1$},
\[
\tilde{H}^q(F_x;\mathbb{Z}) = 0 \quad \text{for all } 0\le q\le n+1,
\]
while satisfying $\chi(F_x)=p\,\chi(F_x/C_p)\equiv 0\pmod{p}$ by placing its
compensating Betti numbers purely in high degrees $q = d_i \gg n+1$.
\end{enumerate}
\end{theorem}

\begin{proof}
Items (i) and (ii) follow immediately from the characterization of compact ENRs
as finite-dimensional ANRs (Borsuk) and the definition of OpenAI's category
$\mathcal{T}(S^d)$ as the idempotent completion of pushforwards from compact
finite polyhedra.  For (iii), take $F_x=S^{2m_i-1}$ with the standard free linear
$C_p$-action and $2m_i-1 > n+1$: then $\tilde{H}^q(F_x;\mathbb{Z})=0$ for all
$q\le n+1$ and $\chi(F_x)=0\equiv 0\pmod{p}$.  (Alternatively, in even top
degree $2m_i>n+1$, a free $C_p$-complex with $\tilde{H}^q=0$ for $q<2m_i$ and
$\tilde{H}^{2m_i}(F_x;\mathbb{Z})\cong\mathbb{Z}^{p-1}$ has
$\chi(F_x)=1+(p-1)=p\equiv 0\pmod{p}$.)  By the Leray spectral sequence of
$\pi_i:E_i\to E_{i-1}$, vanishing of $\tilde{H}^q(F_x;\mathbb{Z})$ for
$q\le n+1$ implies that $\pi_i^*$ is an isomorphism on local cohomology in all
degrees $\le n$, so splicing $F_x$ does \emph{not} destroy local
$n$-manifold homology at stage $i$.
\end{proof}

\subsection{Architecture and exact hurdle for a Dranishnikov--Raymond--Williams construction}
\label{sec:drw-blueprint}

Theorem~\ref{thm:dranishnikov-escape} shows how to avoid both the necklace
pinch points of HM-2 and the low-dimensional homology destruction of
Raymond--Williams (Case~3 of Theorem~\ref{thm:G}): push the free-$C_p$ fiber
topology to growing dimensions $d_i\to\infty$, and kill those high-dimensional
classes at later stages so that $\check{H}^q(X;\mathbb{Z})=\varinjlim_i H^q(E_i;\mathbb{Z})=0$
for all $q>n$ (\emph{pro-acyclicity}).

\begin{proposition}[Pro-acyclicity over $\mathbb{F}_p$ vs $\mathbb{Z}$ in a growing-dimension $C_{p^i}$-tower]
\label{prop:drw-blueprint}
\status{new construction analysis}
Consider an inverse system $E_0\xleftarrow{\pi_1}E_1\xleftarrow{\pi_2}\cdots$
with compatible $C_{p^i}$-actions and growing fiber dimensions
$2m_i-1 > n+1$.
\begin{enumerate}[label=(\alph*)]
\item \textbf{Mod-$p$ pro-acyclicity is automatic ($\mathrm{cd}_p(\mathbb{Z}_p)=1$):}
Unlike the finite group $C_p$, which has $\mathrm{cd}_p(C_p)=\infty$, the
profinite group $\mathbb{Z}_p=\varprojlim C_{p^i}$ has continuous mod-$p$
cohomological dimension $1$: under the quotient $q_{i+1}:C_{p^{i+1}}\twoheadrightarrow C_{p^i}$,
the degree-$2$ Bockstein generator $u_i\in H^2(BC_{p^i};\mathbb{F}_p)$ pulls
back to $q_{i+1}^*(u_i)=0\in H^2(BC_{p^{i+1}};\mathbb{F}_p)$, so
\[
\varinjlim_i H^q(BC_{p^i};\mathbb{F}_p) = H^q_{\mathrm{cts}}(\mathbb{Z}_p;\mathbb{F}_p) = 0
\quad \text{for all } q\ge 2.
\]
Concretely, on the free $C_{p^i}$-sphere $S^{2m_i-1}\subset\mathbb{C}^{m_i}$,
the $q_{i+1}$-equivariant coordinate $p$-th power map
\[
\phi_i(z_1,\dots,z_{m_i}) = \frac{(z_1^p,\dots,z_{m_i}^p)}{\|(z_1^p,\dots,z_{m_i}^p)\|}
\]
has degree $p^{m_i}\equiv 0\pmod{p}$, so $\phi_i^*=0$ on
$\tilde{H}^*(S^{2m_i-1};\mathbb{F}_p)$ in a single step!
\item \textbf{The integral solenoid tail and Dranishnikov killing:}
On integral cohomology $\tilde{H}^{2m_i-1}(S^{2m_i-1};\mathbb{Z})\cong\mathbb{Z}$,
the map $\phi_i^*$ is multiplication by $p^{m_i}\neq 0$.  Iterating $\phi_i$ in
a \emph{fixed} dimension $2m-1$ produces a $p$-adic solenoid with
$\check{H}^{2m-1}\cong\mathbb{Z}[1/p]$ rather than $\dim_{\mathbb{Z}}X=n$.
To obtain $\dim_{\mathbb{Z}}X=n$, stage $i+1$ must instead kill the integral
class in $H^{2m_i-1}(E_i;\mathbb{Z})$ by routing $\pi_{i+1}$ through a higher-dimensional
equivariant Eilenberg--MacLane $K(\mathbb{Z},2m_i-1)$ / Postnikov block of
dimension $d_{i+1}>2m_i-1$ (the Edwards--Walsh / Dranishnikov step).
\item \textbf{The Smith--Coppola constraint in the limit:}
If the limit $X=\varprojlim E_i$ achieves $\dim_{\mathbb{Z}}X=n$, then by
Universal Coefficients $\dim_{\mathbb{F}_p}X\le n<\infty$.  Because P.~A.\
Smith's fixed-point theorem holds for compact spaces of finite
$\mathbb{F}_p$-cohomological dimension (even when $\dim X=\infty$), no
\emph{finite} subgroup $C_p$ can act freely on an $\mathbb{F}_p$-acyclic closed
subset of $X$.  The construction relies on the fact that $\mathbb{Z}_p$ is
\emph{torsion-free} (so no finite $C_p$ acts on $X$ itself) while the orbit
space $X/\mathbb{Z}_p$ absorbs my 1980 transfer / Yang torsion jump
$\dim_{\mathbb{Z}}(X/\mathbb{Z}_p)=n+2$ with top group $\mathbb{Z}[1/p]/\mathbb{Z}$.
\end{enumerate}
\end{proposition}

\begin{checkboxenv}
\textbf{Workplan for a non-necklace HM-6 construction:}
To turn Proposition~\ref{prop:drw-blueprint} into an explicit counterexample in
class HM-6 ($\dim_{\mathbb{Z}}X=n$, $\dim X=\infty$), one must construct the
equivariant Dranishnikov resolution block at step~(b): given the free
$C_{p^i}$-fiber $S^{2m_i-1}$ at stage $i$, build the stage-$(i+1)$ free
$C_{p^{i+1}}$-complex $E_{i+1}\xrightarrow{\pi_{i+1}}E_i$ so that
$\pi_{i+1}^*:H^{2m_i-1}(E_i;\mathbb{Z})\to H^{2m_i-1}(E_{i+1};\mathbb{Z})$ is
zero while maintaining $\tilde{H}^q(F_x;\mathbb{Z})=0$ for all $q\le n+1$.
The subsections below record the Williams obstruction, the Edwards--Walsh
bypass, the Borel-join architecture, and the remaining analytic gate.
\end{checkboxenv}

\subsection{Williams's functor and the 1967 Part~II torus obstruction}
\label{sec:williams-torus}

Williams~\cite{williams1967} (building on Raymond--Williams~\cite{raymondwilliams1960,raymondwilliams1963})
defines a splicing functor $K\Delta_q X$: replace every $m$-simplex of a
barycentric subdivision $K'$ by a building block $(X_i,B_i,q_i)$ carrying a
period-$p^i$ map, glued along $\partial\Delta^m$.  In Part~II
(\cite{williams1967}, pp.~149--154), forcing the bonding map
$L_{i+1}\to L_i$ to factor simplicially through the $n$-skeleton
(Definition~3.1.1: so that $\dim X\le n$) produces a torus
$T=S^1\times\beta^+$ on which a map of type $(1,p)$ cannot respect the
diagonal $C_p$-identification $y_2 y_3^i$.

Concretely, the diagonal circle $S'=\{(\theta,i\theta)\}$ is invariant, and a
map $g'$ forced into the $2$-skeleton has bidegree $(1,p)$ on $T$, hence type
$(i,p)$ on $S'$.  When $i\not\equiv 0\pmod{p}$ this cannot factor through the
diagonal $C_p$-quotient---the classical Part~II no-go.

\begin{remark}[Why Part~II does not kill HM-6]
\label{rmk:williams-depends-on-skeleton}
\status{new}
The obstruction uses the $n$-skeleton retraction hypothesis.  If bonding maps
factor instead through higher cells of an Edwards--Walsh complex
$EW_d(L_i,\mathbb{Z},n)$ that kill $H^{>n}$ \emph{cohomologically} without
collapsing the $(n+1)$-skeleton geometrically, then $S'$ can bound in the
attached cells of dimension $>n$ and the bidegree contradiction disappears.
\end{remark}

\subsection{Edwards--Walsh resolutions and Dranishnikov $K$-theory blocks}
\label{sec:ew-blocks}

The Edwards--Walsh complex $EW(K,\mathbb{Z},n)\to|K|$
\cite{walsh1981,dranishnikov1988,dydakwalsh1993} replaces every simplex
$\sigma$ of dimension $>n$ by an Eilenberg--MacLane complex
$K(\bigoplus\mathbb{Z},n)$ whose $(n+1)$-skeleton is $\sigma^{(n)}$.  This
kills all cohomology in degrees $>n$ (so $\dim_{\mathbb{Z}}X\le n$) while
allowing topological dimension $\dim K_i=d_i\to\infty$.

By Dranishnikov--Ferry--Weinberger~\cite{dfw2006},
$\widetilde{KO}_*(K(\pi,n);\mathbb{Z}_p)=0$ for $n\ge 3$ (and for finite
$\pi$, $n\ge 2$).  Bonding maps may therefore factor through finite skeletons
of $EW(K_i,\mathbb{Z},n)$ while remaining essential in mod-$p$ $K$-theory,
forcing every neighborhood in $X=\varprojlim K_i$ to have infinite covering
dimension.  Dranishnikov's polynomial criterion
\cite[Prop.~3.3]{dranishnikov2016} makes this checkable:
\[
\bigl((y+1)^{p^{k-l}}-1\bigr)^m
\notin
\bigl\langle (y+1)^{p^k}-1,\, y^{n_i}\bigr\rangle
\subset\mathbb{Z}[y]
\]
whenever $m<p^l$ and $m\,p^{k-l}<n_i$, yielding unbounded cup-length on every
open neighborhood $U\subset X$.

\subsection{Borel-join equivariant Edwards--Walsh splicing}
\label{sec:borel-join}

A finite group $C_p$ cannot act freely on an infinite-dimensional
$K(\mathbb{Z},n)$ without nonvanishing Tate cohomology.  The remedy is to
perform Edwards--Walsh attachment via the \textbf{Borel construction} with a
free odd sphere of growing dimension.

\begin{proposition}[Borel-join HM-6 architecture]
\label{prop:borel-join-hm6}
\status{new construction analysis}
Fix $n\ge 2$ and an odd prime $p$.  Inductively construct finite simplicial
complexes $E_i$ with free $C_{p^i}$-actions and bonding maps
$\pi_{i+1}:E_{i+1}\to E_i$ by
\[
E_{i+1}
\;=\;
\bigl(E_i\,\Delta_{q_i}\,EW_{d_i}(E_i,\mathbb{Z},n)\bigr)
\times_{C_{p^{i+1}}}
S^{2m_{i+1}-1},
\]
where $2m_{i+1}-1>n+1$ and $d_{i+1}>d_i$.  Then at every finite stage:
\begin{enumerate}[label=(\roman*)]
\item \textbf{Free action without Smith below degree $n+1$:}
$S^{2m_{i+1}-1}$ is $(2m_{i+1}-2)$-connected with $\chi=0$, so the Borel join
is $(n+1)$-acyclic and $\chi\equiv 0\pmod{p}$ (no Euler contradiction).
\item \textbf{Williams Part~II bypass:}
Edwards--Walsh $(n+2)$-cells let the diagonal circle $S'$ bound without
$n$-skeleton retraction (Remark~\ref{rmk:williams-depends-on-skeleton}).
\item \textbf{Pro-acyclicity $\Rightarrow\dim_{\mathbb{Z}}X=n$:}
coordinate $p$-powers kill mod-$p$ classes in high degree; Edwards--Walsh
cells kill integral classes in degrees $>n$, so
$\varinjlim_i H^q(E_i;\mathbb{Z})=0$ for $q>n$.
\item \textbf{Infinite covering dimension on every neighborhood:}
Dranishnikov's criterion supplies cup-length $\to\infty$ on every open star,
hence $\dim U=\infty$ while $\dim_{\mathbb{Z}}U=n$.
\item \textbf{BRW constraint on the orbit space:}
any successful limit must satisfy
$\dim_{\mathbb{Z}}(X/A_p)=n+2$ with top group injecting onto
$R_p(1)=\mathbb{Z}[1/p]/\mathbb{Z}$
\cite{raymondwilliams1960,coppola1980}---a constraint, not an obstruction.
\end{enumerate}
\end{proposition}

\begin{proof}[Proof sketch]
Item~(i) is the standard connectivity of free linear $C_p$-spheres and the
join formula $\mathrm{conn}(F\times_{C_p}S^{2m-1})\ge 2m-2$ when the action on
the sphere is free.  Item~(ii) is Remark~\ref{rmk:williams-depends-on-skeleton}.
Item~(iii) combines Proposition~\ref{prop:drw-blueprint}(a)--(b) with the
Postnikov killing cells of $EW_{d_i}$.  Item~(iv) is
\cite[Prop.~3.3--3.4]{dranishnikov2016}.  Item~(v) is the classical
Bredon--Raymond--Williams transfer theorem, which uses only
$\dim_{\mathbb{Z}}X=n<\infty$, not finite covering dimension.
\end{proof}

\subsection{Remaining analytic gate and computational verification}
\label{sec:hm6-gate}

\begin{theorem}[Finite-stage HM-6 gates]
\label{thm:hm6-gates}
\status{new --- computational / Lean}
On every finite stage of the architecture in
Proposition~\ref{prop:borel-join-hm6}, the following four gates hold for all
tested parameters $(n,p,\text{stages})\in\{(5,3,4),(7,5,3),(6,3,5),(9,7,4)\}$
(symbolic harness
\texttt{scripts/hm6\_equivariant\_ew\_harness.py}; loop log
\texttt{counterexamples/hilbert-smith/HM6\_LOOP\_LOG.md}):
\begin{description}[style=nextline,leftmargin=1.2em]
\item[G1] Williams torus bypass via Edwards--Walsh $(n+2)$-cells.
\item[G2] Pro-acyclicity in degrees $>n$ (integral and mod~$p$) with
strictly growing $\dim E_i$.
\item[G3] Local $\mathbb{Z}$-cohomology $n$-manifold condition
$H^k\cong\delta_{k,n}\mathbb{Z}$, with orbit-space Yang jump $n+2$ and
$R_p(1)$ tower.
\item[G4] Dranishnikov Prop.~3.3 polynomial criterion with unbounded
cup-length $\Rightarrow\dim U=\infty$.
\end{description}
The Borel-join connectivity lemma, the cup-length valuation gate
$v_p(\mathrm{coeff}_m)=m<k$, and the HM-6 escape from
\texttt{OpenAIENRInput}/\texttt{DabrowskiTowerInput} are formalized in
\texttt{lean/CoppolaHS/CoppolaHS/DranishnikovHM6.lean}
(Lean~\texttt{v4.35.0-rc3}, \texttt{lake build} clean).
\end{theorem}

\begin{remark}[Single remaining analytic verification]
\label{rmk:hm6-remaining}
\status{open geometric glue}
What remains for a complete geometric counterexample is the local
$clc_{\mathbb{Z}}$ degree-$n$ Poincar\'e duality \emph{gluing} across simplex
boundaries $\partial\sigma$ when the Edwards--Walsh block is pulled back over
the Yang $(n+2)$-dimensional orbit space.  Finite-stage gates G1--G4 verify
the algebraic and homological scaffolding; the continuous-limit gluing is the
sole outstanding geometric step.
\end{remark}

\begin{checkboxenv}
Confirm G1--G4 on the harness for your preferred $(n,p)$; check that
\texttt{DranishnikovHM6.lean} builds; treat Remark~\ref{rmk:hm6-remaining} as
the remaining research gate, not as a claimed existence theorem.
\end{checkboxenv}

%======================================================================
\section{Summary matrix}
\label{sec:matrix}

\begin{center}
\small
\begin{tabular}{@{}p{3.4cm}p{2.4cm}p{7.2cm}@{}}
\toprule
Class & Faithful $\mathbb{Z}_p$? & Verdict \\
\midrule
HM-1 non-$clc$ / coeff.\ mismatch & yes (mechanism) & Outside both 2026 scopes \\
HM-2 necklace & yes (explicit) & Local homology fails at junction \\
HM-3 polyhedral & no & Thm~\ref{thm:C} $\to$ manifold theorem \\
HM-4 resolvable / $X_{\mathrm{man}}\neq\varnothing$ & no & Thm~\ref{thm:D},~\ref{thm:settled-chart} \\
HM-5 BFMW ENR, $X_{\mathrm{man}}=\varnothing$, $I\neq 1$ & no (rel.\ E or F) & Bounded-dim inverse limit blocked by Thm~\ref{thm:G} \\
HM-6 non-ENR, $\dim_{\mathbb{Z}}X=n$, $\dim X=\infty$ & \textbf{open candidate} & Escapes E, F, G(2); Borel-join architecture Prop.~\ref{prop:borel-join-hm6}; gates Thm~\ref{thm:hm6-gates}; glue Rem.~\ref{rmk:hm6-remaining} \\
\bottomrule
\end{tabular}
\end{center}

\noindent\textit{Reading guide for HM-5 vs.\ HM-6.}
Within finite topological dimension ($\dim X=n$, i.e.\ ENRs, HM-5),
exclusion holds relative to Theorem~\ref{thm:E} or~\ref{thm:F}, and
Theorem~\ref{thm:G} blocks bounded-dimension inverse limits.  Dropping finite
topological dimension while keeping finite cohomological dimension
($\dim_{\mathbb{Z}}X=n$, $\dim X=\infty$, HM-6) is the unique frontier where a
non-necklace inverse-limit construction is not ruled out by any of the three.

%======================================================================
\section*{Acknowledgments}
I thank Louis~F.\ McAuley for the 1979--80 Binghamton apprenticeship that
made the transfer machine possible.  The 2026 OpenAI and Dabrowski manuscripts
are the occasion for this synthesis; any errors in delimitation are mine.
Cursor(Opus~5.5/Grok) assisted with drafting and typesetting under my
direction.

\appendix
%======================================================================
\section{Lean status (honest boundary)}
\label{app:lean}

Dabrowski's formalization~\cite{dabrowski2026} was compiled locally from branch
\texttt{palomar-review/module-system-20261015} at
\texttt{b75523f0036bc857f31620ee18af970151036311}
(vendored under \texttt{dabrowski/upstream-palomar/}), Lean
\texttt{v4.35.0-rc3}, mathlib \texttt{1a547d8}, TauCeti \texttt{c7af81f}.

\begin{verbatim}
#print axioms HSFormal.hilbertSmith
'HSFormal.hilbertSmith' depends on axioms:
  [propext, Classical.choice, Quot.sound]
\end{verbatim}

Extension modules live in \texttt{lean/CoppolaHS/}:
\begin{itemize}
\item \texttt{CyclotomicTower.lean} --- product of $k$ prime-power cyclotomics
$\Rightarrow p^k\mid F(1)$ (Theorem~\ref{thm:F} arithmetic);
\item \texttt{QuinnIndex.lean} --- abstract $I(X):\mathbb{Z}$ with
$I\equiv 1\pmod{8}$ and product-formula hypotheses;
\item \texttt{PadicExclusionENR.lean} --- conditional Theorems E/F shells;
\item \texttt{ChallengeENR.lean} / \texttt{SolutionENR.lean} --- statement +
relative solution notes;
\item \texttt{DranishnikovHM6.lean} --- Borel-join connectivity, Dranishnikov
cup-length valuation gate, HM-6 escape from ENR packaging
(Theorem~\ref{thm:hm6-gates}).
\end{itemize}
We do \textbf{not} claim a full controlled $L$-theory reimplementation for ENRs
in Lean~v1.  Geometric BFMW inputs remain named hypotheses with paper citations.
The HM-6 geometric gluing of Remark~\ref{rmk:hm6-remaining} is likewise not
formalized.

%======================================================================
\section{Reviewer checklist}
\label{app:checklist}

Mark each item PASS / WEAK / FAIL.  Co-author invitations may attach initials.

\vspace{0.6em}
\noindent
\begin{tabular}{@{}clp{7.8cm}@{}}
\toprule
 & Statement & Your mark \\
\midrule
\checkbox & Lemma~\ref{lem:A} (transfer over $\mathbb{Q}$) & \underline{\hspace{2.5cm}} \\
\checkbox & Lemma~\ref{lem:B} (Yang $p$-torsion jump) & \underline{\hspace{2.5cm}} \\
\checkbox & Corollary~\ref{cor:mcauley} (McAuley break) & \underline{\hspace{2.5cm}} \\
\checkbox & Proposition~\ref{prop:necklace} (HM-2 necklace) & \underline{\hspace{2.5cm}} \\
\checkbox & Theorem~\ref{thm:C} (polyhedral reduction) & \underline{\hspace{2.5cm}} \\
\checkbox & Theorem~\ref{thm:D} (resolvable / nonempty stratum) & \underline{\hspace{2.5cm}} \\
\checkbox & Theorem~\ref{thm:dabrowski-gap} (single-$p$ gap) & \underline{\hspace{2.5cm}} \\
\checkbox & Theorem~\ref{thm:E} (OpenAI ENR rephrase) & \underline{\hspace{2.5cm}} \\
\checkbox & Theorem~\ref{thm:F} ($k$-tower upgrade) & \underline{\hspace{2.5cm}} \\
\checkbox & Prop.~\ref{prop:short-answers} (settled vs open) & \underline{\hspace{2.5cm}} \\
\checkbox & Theorem~\ref{thm:settled-chart} /~\ref{thm:no-ilim-enr} & \underline{\hspace{2.5cm}} \\
\checkbox & Theorem~\ref{thm:G} (bounded-dim inverse-limit trichotomy) & \underline{\hspace{2.5cm}} \\
\checkbox & Corollary~\ref{cor:H} (necklace escape in bounded dim) & \underline{\hspace{2.5cm}} \\
\checkbox & Theorem~\ref{thm:dranishnikov-escape} ($\dim_{\mathbb{Z}}=n,\dim=\infty$ escape) & \underline{\hspace{2.5cm}} \\
\checkbox & Prop.~\ref{prop:drw-blueprint} (Dranishnikov--RW blueprint) & \underline{\hspace{2.5cm}} \\
\checkbox & Prop.~\ref{prop:borel-join-hm6} (Borel-join EW + Williams) & \underline{\hspace{2.5cm}} \\
\checkbox & Theorem~\ref{thm:hm6-gates} (G1--G4 harness / Lean) & \underline{\hspace{2.5cm}} \\
\checkbox & Remark~\ref{rmk:hm6-remaining} (remaining $clc_{\mathbb{Z}}$ glue) & \underline{\hspace{2.5cm}} \\
\bottomrule
\end{tabular}

\vspace{1em}
\noindent\textit{Priority for first readers:} Section~\ref{sec:plain-english}
(plain-English guide to Theorems~\ref{thm:dabrowski-gap}--\ref{thm:F}),
Section~\ref{sec:dranishnikov}--\ref{sec:drw-blueprint} (the
$\dim_{\mathbb{Z}}X=n$, $\dim X=\infty$ frontier for non-necklace examples),
then Theorem~\ref{thm:G} and Corollary~\ref{cor:mcauley}.

\begin{thebibliography}{99}

\bibitem{bfmw1996}
J.~Bryant, S.~Ferry, W.~Mio, and S.~Weinberger,
\textit{Topology of homology manifolds},
Ann.\ of Math.\ (2) \textbf{143} (1996), 435--467.

\bibitem{dranishnikov1988}
A.~N.\ Dranishnikov,
\textit{On a problem of P.~S.\ Aleksandrov},
Mat.\ Sb.\ (N.S.) \textbf{135(177)} (1988), 551--557; English transl.,
Math.\ USSR-Sb.\ \textbf{63} (1989), 539--545.

\bibitem{dydakwalsh1993}
J.~Dydak and J.~J.\ Walsh,
\textit{Infinite-dimensional compacta having cohomological dimension two: an
application of the Sullivan conjecture},
Topology \textbf{32} (1993), 93--104.

\bibitem{dfw2006}
A.~N.\ Dranishnikov, S.~Ferry, and S.~Weinberger,
\textit{Large Riemannian manifolds which are flexible},
Ann.\ of Math.\ (2) \textbf{157} (2003), 919--938;
see also arXiv:math/0611004.

\bibitem{dranishnikov2016}
A.~N.\ Dranishnikov,
\textit{On dimension of product of groups},
Sbornik Math.\ \textbf{207} (2016); see also arXiv related notes on
essential lens sequences and $K$-theory cup-length.

\bibitem{williams1967}
R.~F.\ Williams,
\textit{The construction of certain $p$-adic transformation groups},
Trans.\ Amer.\ Math.\ Soc.\ \textbf{129} (1967), 140--156.

\bibitem{mitchell1990}
W.~J.~R.\ Mitchell,
\textit{Defining the boundary of a homology manifold},
Proc.\ Amer.\ Math.\ Soc.\ \textbf{110} (1990), 509--513.

\bibitem{walsh1981}
J.~J.\ Walsh,
\textit{Dimension, cohomological dimension, and cell-like mappings},
in \textit{Shape Theory and Geometric Topology} (Dubrovnik, 1981),
Lecture Notes in Math.\ \textbf{870}, Springer, 1981, pp.~105--118.

\bibitem{connerfloyd1959}
P.~E.\ Conner and E.~E.\ Floyd,
\textit{Differentiable periodic maps},
Springer, 1960.  %(classic reference; Newman-type results)

\bibitem{coppola1980}
A.~J.\ Coppola,
\textit{On $p$-adic transformation groups},
Ph.D.\ thesis, SUNY at Binghamton, 1980
(summarized in McAuley~\cite{mcauley2001}, ref.~[28]).

\bibitem{dabrowski2026}
A.~Dabrowski (with AI agents),
\textit{Integral tail signatures and the Hilbert--Smith conjecture},
preprint, 8~Oct~2026,
\url{https://github.com/adbrw/HS_proof}.

\bibitem{mcauley2001}
L.~F.\ McAuley,
\textit{A proof of the Hilbert--Smith conjecture},
arXiv:math/0103215v8, 2001.

\bibitem{newman1931}
M.~H.~A.\ Newman,
\textit{A theorem on periodic transformations of spaces},
Quart.\ J.\ Math.\ Oxford Ser.\ \textbf{2} (1931), 1--8.

\bibitem{openai2026}
OpenAI,
\textit{The Hilbert--Smith conjecture in every finite dimension},
preprint, 23~Sep~2026,
\url{https://github.com/openai/math}.

\bibitem{raymond1965}
F.~Raymond,
\textit{Some remarks on the coefficients used in the theory of homology manifolds},
Pacific J.\ Math.\ \textbf{15} (1965), 1365--1376.

\bibitem{raymondwilliams1960}
F.~Raymond and R.~F.\ Williams,
\textit{Examples of $p$-adic transformation groups},
Bull.\ Amer.\ Math.\ Soc.\ \textbf{66} (1960), 392--394.

\bibitem{raymondwilliams1963}
F.~Raymond and R.~F.\ Williams,
\textit{Examples of $p$-adic transformation groups},
Ann.\ of Math.\ (2) \textbf{78} (1963), 92--106.

\bibitem{smith1941}
P.~A.\ Smith,
\textit{Transformations of finite period},
Ann.\ of Math.\ (2) \textbf{39} (1938), 127--164; III,
ibid.\ \textbf{42} (1941), 446--458.

\bibitem{yang1960}
C.~T.\ Yang,
\textit{$p$-adic transformation groups},
Michigan Math.\ J.\ \textbf{7} (1960), 201--218.

\end{thebibliography}

\end{document}
